\section{Statement, conventions and proof structure}\label{sec:intro}
A digraph is loopless and has at most one arc per ordered pair; opposite arcs are allowed. A proper colouring has distinct colours on arcs leaving each vertex and on arcs entering each vertex. An incoming and an outgoing arc at the same vertex may have the same colour. All paths below have distinct physical vertices; their length is their number of arcs.
\begin{theorem}[Fixed positive density]\label{gen:theorem}
For every fixed $\alpha>0$, all sufficiently large loopless simple digraphs on $n$ vertices with indegree and outdegree $d\ge\alpha{}n$, properly coloured by exactly $d$ colours, contain a rainbow directed path with $d-1$ arcs.
\end{theorem}
\begin{corollary}[The undirected conclusion]\label{gen:undirected}
For every fixed $\alpha>0$, all sufficiently large simple $d$-regular graphs on $n$ vertices with $d\ge\alpha{}n$, properly coloured by exactly $d$ colours, contain a rainbow path of length $d-1$.
\end{corollary}
\begin{proof}
Replace each edge $xy$ by $x\to{}y$ and $y\to{}x$, both carrying its colour. The resulting digraph satisfies Theorem~\ref{gen:theorem}. Its rainbow directed simple path gives the required undirected simple path.
\end{proof}

Every colour in the theorem is a directed cycle cover, written $x\mapsto{}P_c(x)$. Its incidence graph has copies $V^+,V^-$ and colour-$c$ matching $x^+P_c(x)^-$. Write $I_c$ for the matching involution and $\iota(x^+)=x^-$ for physical identification. A set of copies is \emph{admissible} if it uses at most one copy of each physical vertex. In rectangular restrictions, identification is a partial bijection between the two classes. No real edge joins copies with the same physical label. Incidence walks are counting devices; only forward arcs and the final directed assembly represent directed paths.

All constants are fixed before the ambient order tends to infinity. The notation $x\ll{}y$ means that $x$ is chosen sufficiently small after $y$; $o(m)$ estimates are uniform over the vertices, colours and finitely many blocks concerned. Forbidden vertices always mean their entire physical resource, including a mate when present. A prescribed root is exempt only where explicitly stated. Artificial edges used to complete partial permutations are counting devices and never output resources.

We use the bounded-degree distributive template and weighted graph decomposition/connection estimates of \cite[Section 3.9]{Montgomery}, the compatible even-cycle input in \cite[Theorem 3.4]{BMMP}, and the linear-hypergraph edge-colouring estimate of \cite{MolloyReed}. Their applications, including the quantitative rounding loss, are stated at the points of use. The proof first constructs robust local absorption and exact matching modules, then routes and balances approximate spectral blocks. A decomposition of the original \emph{real} incidence matrix supplies those blocks. Finally a short prefix uses exceptional colours and enforces the exact physical capacity needed for a suffix of the claimed length.
